% GENERATED FILE. Edit canonical lessons, then run npm run generate:books.
% canonical-source-hash: fnv1a64:47a5dd542e58d90f
% canonical-lessons: UR-C439-gidar, UR-C439-jhingur, UR-C439-maveshi, UR-C439-charna, UR-C439-jotna, UR-R439-review-433-439, UR-R439-review-final-439

\chapter{Night in the Fields}
\label{ch:ur-night-in-the-fields}
\hlchaptermodality{439}

\begin{chapteropening}
\textbf{By the end of this chapter:} \emph{I can describe the fields at night and at work: jackals, crickets, grazing cattle and the plough.}

Complete the last lesson of chapter 439: I can describe the fields at night and at work: jackals, crickets, grazing cattle and the plough.
\end{chapteropening}

\section[\texorpdfstring{\ur{گیدڑ} (gīdaṛ)}{گیدڑ (gīdaṛ)}]{\ur{گیدڑ} (gīdaṛ) — a jackal}

\label{lesson:UR-C439-gidar}



\begin{quote}
Before the new one: say the Urdu for to peck (at), then the Urdu for prey.
\end{quote}

\subsection*{You'll want to know: \ur{گیدڑ}}
\textbf{\ur{گیدڑ}} — \emph{gīdaṛ} — \textquotedblleft{}a jackal\textquotedblright{}.

\begin{quote}
\ur{رات} \ur{کو} \ur{کھیتوں} \ur{میں} \ur{گیدڑ} \ur{بولتے} \ur{ہیں۔}
\end{quote}

\begin{quote}
\emph{rāt ko khetoṅ meṅ gīdaṛ bolte haiṅ.}
\end{quote}

\begin{quote}
At night jackals howl in the fields.
\end{quote}

Masculine.

\begin{cousinweb}
\textbf{\ur{گیدڑ}} (\emph{gīdaṛ}) is inherited: CDIAL lists the Hindi cousin \emph{gīdaṛ}, \textquotedblleft{}a jackal\textquotedblright{}, under \emph{gidda}, a form it reconstructs (CDIAL 4158).
\end{cousinweb}

\begin{grammarlens}[title={What you've built}]
One more word for getting around.
\end{grammarlens}

\subsection*{Your turn}
\begin{itemize}
\raggedright
  \item \emph{Say it:} \emph{gīdaṛ}
  \item \emph{Say it:} \emph{gīdaṛ}, once more
  \item \emph{Say it:} \emph{rāt ko khetoṅ meṅ gīdaṛ bolte haiṅ.} — At night jackals howl in the fields.
  \item \emph{Recall:} say the Urdu for to peck (at), then the Urdu for prey, then say \emph{gīdaṛ} again
\end{itemize}

\begin{tcolorbox}[breakable,colback=teal!4,colframe=teal!35!black,title={Before you move on}]
What does \ur{گیدڑ} mean? (\textquotedblleft{}A jackal\textquotedblright{}.) Say it once more.
\end{tcolorbox}
\section[\texorpdfstring{\ur{جھینگر} (jhīṅgur)}{جھینگر (jhīṅgur)}]{\ur{جھینگر} (jhīṅgur) — a cricket (insect)}

\label{lesson:UR-C439-jhingur}



\begin{quote}
Before the new one: say the Urdu for prey, then the Urdu for a jackal.
\end{quote}

\subsection*{You'll want to know: \ur{جھینگر}}
\textbf{\ur{جھینگر}} — \emph{jhīṅgur} — \textquotedblleft{}a cricket (insect)\textquotedblright{}.

\begin{quote}
\ur{رات} \ur{کو} \ur{جھینگر} \ur{شور} \ur{کرتے} \ur{ہیں۔}
\end{quote}

\begin{quote}
\emph{rāt ko jhīṅgur shor karte haiṅ.}
\end{quote}

\begin{quote}
At night the crickets chirp.
\end{quote}

Masculine.

\begin{cousinweb}
\textbf{\ur{جھینگر}} (\emph{jhīṅgur}) is inherited: CDIAL lists the Hindi cousin \emph{jhĩgur}, \textquotedblleft{}a cricket\textquotedblright{}, under \emph{jhiṅgura}, a form it reconstructs (CDIAL 5387).
\end{cousinweb}

\begin{grammarlens}[title={What you've built}]
One more word for getting around.
\end{grammarlens}

\subsection*{Your turn}
\begin{itemize}
\raggedright
  \item \emph{Say it:} \emph{jhīṅgur}
  \item \emph{Say it:} \emph{jhīṅgur}, once more
  \item \emph{Say it:} \emph{rāt ko jhīṅgur shor karte haiṅ.} — At night the crickets chirp.
  \item \emph{Recall:} say the Urdu for prey, then the Urdu for a jackal, then say \emph{jhīṅgur} again
\end{itemize}

\begin{tcolorbox}[breakable,colback=teal!4,colframe=teal!35!black,title={Before you move on}]
What does \ur{جھینگر} mean? (\textquotedblleft{}A cricket (insect)\textquotedblright{}.) Say it once more.
\end{tcolorbox}
\section[\texorpdfstring{\ur{مویشی} (maveshī)}{مویشی (maveshī)}]{\ur{مویشی} (maveshī) — livestock; cattle}

\label{lesson:UR-C439-maveshi}



\begin{quote}
Before the new one: say the Urdu for a jackal, then the Urdu for a cricket (insect).
\end{quote}

\subsection*{You'll want to know: \ur{مویشی}}
\textbf{\ur{مویشی}} — \emph{maveshī} — \textquotedblleft{}livestock; cattle\textquotedblright{}.

\begin{quote}
\ur{کسان} \ur{کے} \ur{مویشی} \ur{کھیت} \ur{میں} \ur{ہیں۔}
\end{quote}

\begin{quote}
\emph{kisān ke maveshī khet meṅ haiṅ.}
\end{quote}

\begin{quote}
The farmer's livestock are in the field.
\end{quote}

Masculine, used in the plural.

\begin{cousinweb}
\textbf{\ur{مویشی}} (\emph{maveshī}) is usually traced, through Persian, to Arabic \emph{mawāshī}, \textquotedblleft{}livestock, herds\textquotedblright{}.
\end{cousinweb}

\begin{grammarlens}[title={What you've built}]
One more word for getting around.
\end{grammarlens}

\subsection*{Your turn}
\begin{itemize}
\raggedright
  \item \emph{Say it:} \emph{maveshī}
  \item \emph{Say it:} \emph{maveshī}, once more
  \item \emph{Say it:} \emph{kisān ke maveshī khet meṅ haiṅ.} — The farmer's livestock are in the field.
  \item \emph{Recall:} say the Urdu for a jackal, then the Urdu for a cricket (insect), then say \emph{maveshī} again
\end{itemize}

\begin{tcolorbox}[breakable,colback=teal!4,colframe=teal!35!black,title={Before you move on}]
What does \ur{مویشی} mean? (\textquotedblleft{}Livestock; cattle\textquotedblright{}.) Say it once more.
\end{tcolorbox}
\section[\texorpdfstring{\ur{چرنا} (charnā)}{چرنا (charnā)}]{\ur{چرنا} (charnā) — to graze}

\label{lesson:UR-C439-charna}



\begin{quote}
Before the new one: say the Urdu for a cricket (insect), then the Urdu for livestock.
\end{quote}

\subsection*{You'll want to know: \ur{چرنا}}
\textbf{\ur{چرنا}} — \emph{charnā} — \textquotedblleft{}to graze\textquotedblright{}.

\begin{quote}
\ur{مویشی} \ur{دن} \ur{بھر} \ur{میدان} \ur{میں} \ur{چرتے} \ur{ہیں۔}
\end{quote}

\begin{quote}
\emph{maveshī din bhar maidān meṅ charte haiṅ.}
\end{quote}

\begin{quote}
The cattle graze in the meadow all day.
\end{quote}

Past (it …): \ur{وہ} \ur{گائے} \ur{دیر} \ur{تک} \ur{چری}, \emph{voh gāe der tak charī}.

\begin{cousinweb}
\textbf{\ur{چرنا}} (\emph{charnā}) is inherited: CDIAL lists the Hindi cousin \emph{carnā}, \textquotedblleft{}to move about, to graze\textquotedblright{}, under Sanskrit \emph{cárati}, \textquotedblleft{}moves\textquotedblright{} (CDIAL 4686).
\end{cousinweb}

\begin{grammarlens}[title={What you've built}]
One more word for getting around.
\end{grammarlens}

\subsection*{Your turn}
\begin{itemize}
\raggedright
  \item \emph{Say it:} \emph{charnā}
  \item \emph{Say it:} \emph{charnā}, once more
  \item \emph{Say it:} \emph{maveshī din bhar maidān meṅ charte haiṅ.} — The cattle graze in the meadow all day.
  \item \emph{Recall:} say the Urdu for a cricket (insect), then the Urdu for livestock, then say \emph{charnā} again
\end{itemize}

\begin{tcolorbox}[breakable,colback=teal!4,colframe=teal!35!black,title={Before you move on}]
What does \ur{چرنا} mean? (\textquotedblleft{}To graze\textquotedblright{}.) Say it once more.
\end{tcolorbox}
\section[\texorpdfstring{\ur{جوتنا} (jotnā)}{جوتنا (jotnā)}]{\ur{جوتنا} (jotnā) — to plough; to yoke}

\label{lesson:UR-C439-jotna}



\begin{quote}
Before the new one: say the Urdu for livestock, then the Urdu for to graze.
\end{quote}

\subsection*{You'll want to know: \ur{جوتنا}}
\textbf{\ur{جوتنا}} — \emph{jotnā} — \textquotedblleft{}to plough; to yoke\textquotedblright{}.

\begin{quote}
\ur{کسان} \ur{صبح} \ur{سویرے} \ur{کھیت} \ur{جوتتا} \ur{ہے۔}
\end{quote}

\begin{quote}
\emph{kisān subah savere khet jottā hai.}
\end{quote}

\begin{quote}
The farmer ploughs his field early in the morning.
\end{quote}

Past, with \ur{نے} (he or she …): \ur{اس} \ur{نے} \ur{دو} \ur{کھیت} \ur{جوتے}, \emph{us ne do khet jote}.

\begin{cousinweb}
\textbf{\ur{جوتنا}} (\emph{jotnā}) is inherited: CDIAL lists the Hindi cousin \emph{jotnā}, \textquotedblleft{}to yoke, to plough\textquotedblright{}, under Sanskrit \emph{yoktrayati}, \textquotedblleft{}fastens\textquotedblright{} (CDIAL 10524).
\end{cousinweb}

\begin{grammarlens}[title={What you've built}]
One more word for getting around.
\end{grammarlens}

\subsection*{Your turn}
\begin{itemize}
\raggedright
  \item \emph{Say it:} \emph{jotnā}
  \item \emph{Say it:} \emph{jotnā}, once more
  \item \emph{Say it:} \emph{kisān subah savere khet jottā hai.} — The farmer ploughs his field early in the morning.
  \item \emph{Recall:} say the Urdu for livestock, then the Urdu for to graze, then say \emph{jotnā} again
\end{itemize}

\begin{tcolorbox}[breakable,colback=teal!4,colframe=teal!35!black,title={Before you move on}]
What does \ur{جوتنا} mean? (\textquotedblleft{}To plough; to yoke\textquotedblright{}.) Say it once more.
\end{tcolorbox}
\section[(dialogue)]{Looking back: from Mounds, Heaps and Caves to Night in the Fields}

\label{lesson:UR-R439-review-433-439}



\begin{quote}
Say the Urdu for a mound and a heap.
\end{quote}

\subsection*{Your turn}
Say these aloud:

\begin{itemize}
\raggedright
  \item a mound, a heap, a cave, a straw, to be uprooted
  \item a billow, a marsh, a landing place (on a river), to toss (up), to dip
  \item a bay, an anchor, a net, a seashell, a dip (under water)
  \item nature, a fragrance, a new shoot, a petal, to wilt
  \item a bush, shade, dense, moss, to come up (of plants)
  \item a beak, a wing, a nest, to peck (at), prey
  \item a jackal, a cricket (insect), livestock, to graze, to plough
\end{itemize}

\begin{tcolorbox}[breakable,colback=teal!4,colframe=teal!35!black,title={Before you move on}]
A mound? (\textbf{\ur{ٹیلا}}, \emph{ṭīlā}.) A heap? (\textbf{\ur{ڈھیر}}, \emph{ḍher}.) A cave? (\textbf{\ur{غار}}, \emph{ghār}.) A straw? (\textbf{\ur{تنکا}}, \emph{tinkā}.) To be uprooted? (\textbf{\ur{اکھڑنا}}, \emph{ukhaṛnā}.) A billow? (\textbf{\ur{موج}}, \emph{mauj}.) A marsh? (\textbf{\ur{دلدل}}, \emph{daldal}.) A landing place (on a river)? (\textbf{\ur{گھاٹ}}, \emph{ghāṭ}.) To toss (up)? (\textbf{\ur{اچھالنا}}, \emph{uchhālnā}.) To dip? (\textbf{\ur{ڈبونا}}, \emph{ḍubonā}.) A bay? (\textbf{\ur{خلیج}}, \emph{khalīj}.) An anchor? (\textbf{\ur{لنگر}}, \emph{langar}.) A net? (\textbf{\ur{جال}}, \emph{jāl}.) A seashell? (\textbf{\ur{سیپ}}, \emph{sīp}.) A dip (under water)? (\textbf{\ur{ڈبکی}}, \emph{ḍubkī}.) Nature? (\textbf{\ur{قدرت}}, \emph{qudrat}.) A fragrance? (\textbf{\ur{خوشبو}}, \emph{khushbū}.) A new shoot? (\textbf{\ur{کونپل}}, \emph{koṅpal}.) A petal? (\textbf{\ur{پنکھڑی}}, \emph{pankhuṛī}.) To wilt? (\textbf{\ur{مرجھانا}}, \emph{murjhānā}.) A bush? (\textbf{\ur{جھاڑی}}, \emph{jhāṛī}.) Shade? (\textbf{\ur{سایہ}}, \emph{sāyā}.) Dense? (\textbf{\ur{گھنا}}, \emph{ghanā}.) Moss? (\textbf{\ur{کائی}}, \emph{kāī}.) To come up (of plants)? (\textbf{\ur{اگنا}}, \emph{ugnā}.) A beak? (\textbf{\ur{چونچ}}, \emph{choṅch}.) A wing? (\textbf{\ur{پنکھ}}, \emph{pankh}.) A nest? (\textbf{\ur{آشیانہ}}, \emph{āshiyānā}.) To peck (at)? (\textbf{\ur{چگنا}}, \emph{chugnā}.) Prey? (\textbf{\ur{شکار}}, \emph{shikār}.) A jackal? (\textbf{\ur{گیدڑ}}, \emph{gīdaṛ}.) A cricket (insect)? (\textbf{\ur{جھینگر}}, \emph{jhīṅgur}.) Livestock? (\textbf{\ur{مویشی}}, \emph{maveshī}.) To graze? (\textbf{\ur{چرنا}}, \emph{charnā}.) To plough? (\textbf{\ur{جوتنا}}, \emph{jotnā}.)
\end{tcolorbox}
\section[(dialogue)]{Night in the Fields — once more}

\label{lesson:UR-R439-review-final-439}



\begin{quote}
Say the Urdu for a jackal and a cricket (insect).
\end{quote}

\subsection*{Your turn}
\begin{itemize}
\raggedright
  \item \emph{Say it:} a jackal, a cricket (insect), livestock, to graze, to plough
\end{itemize}

\begin{tcolorbox}[breakable,colback=teal!4,colframe=teal!35!black,title={Before you move on}]
A jackal? (\textbf{\ur{گیدڑ}}, \emph{gīdaṛ}.) A cricket (insect)? (\textbf{\ur{جھینگر}}, \emph{jhīṅgur}.) Livestock? (\textbf{\ur{مویشی}}, \emph{maveshī}.) To graze? (\textbf{\ur{چرنا}}, \emph{charnā}.) To plough? (\textbf{\ur{جوتنا}}, \emph{jotnā}.)
\end{tcolorbox}
